\chapter{The Sisal-2026 Type System}
\label{chap:type_system}

\section{Scalar Primitive Types}
Sisal-2026 provides a comprehensive suite of scalar primitive integer and floating-point types:
\begin{table}[h!]
\centering
\begin{tabularx}{\textwidth}{llX}
\toprule
\textbf{Type Name} & \textbf{C++23 Representation} & \textbf{Description} \\
\midrule
\texttt{integer} & \texttt{int32\_t} & 32-bit signed integer \\
\texttt{long} & \texttt{int64\_t} & 64-bit signed integer \\
\texttt{unsigned} & \texttt{uint32\_t} & 32-bit unsigned integer \\
\texttt{unsigned long} & \texttt{uint64\_t} & 64-bit unsigned integer \\
\texttt{real} & \texttt{float} & Single-precision IEEE 754 float (32-bit) \\
\texttt{double} & \texttt{double} & Double-precision IEEE 754 float (64-bit) \\
\texttt{boolean} & \texttt{bool} & Boolean logical value (\texttt{true}/\texttt{false}) \\
\texttt{character} & \texttt{char} & 8-bit ASCII character \\
\bottomrule
\end{tabularx}
\caption{Primitive Scalar Types in Sisal-2026.}
\label{tab:primitive_types}
\end{table}

\section{Strict Type Safety \& Explicit Conversions}
Sisal-2026 enforces \textbf{strict static type safety}. Binary operators and function calls prohibit implicit type coercion across distinct scalar types. Operands in expressions must match exactly; mixed-type arithmetic (such as adding an \texttt{integer} to a \texttt{double} or a \texttt{real} to an \texttt{unsigned}) triggers a static compilation type error.

To combine different scalar types, operands must be explicitly converted using built-in type conversion operators:

\begin{table}[h!]
\centering
\begin{tabularx}{\textwidth}{lX}
\toprule
\textbf{Conversion Function} & \textbf{Semantics \& Behavior} \\
\midrule
\texttt{double(x)} & Promotes \texttt{integer}, \texttt{long}, or \texttt{real} to double-precision \texttt{double}. \\
\texttt{real(x)} & Converts integer types or demotes \texttt{double} to single-precision \texttt{real}. \\
\texttt{integer(x)} & Truncates floating-point types or demotes 64-bit \texttt{long} to 32-bit \texttt{integer}. \\
\texttt{long(x)} & Promotes 32-bit \texttt{integer} to 64-bit \texttt{long}. \\
\texttt{unsigned(x)} & Converts signed integer to 32-bit \texttt{unsigned}. \\
\texttt{unsigned\_long(x)} & Converts signed/unsigned types to 64-bit \texttt{unsigned long}. \\
\bottomrule
\end{tabularx}
\caption{Explicit Type Conversion Operators in Sisal-2026.}
\label{tab:type_conversions}
\end{table}

\begin{lstlisting}[caption={Explicit Type Conversion Example}]
function MixedArithmetic( i : integer; d : double returns double )
  % ERROR: i + d  (Implicit mixing of integer and double is forbidden!)
  
  % CORRECT: Explicitly promote i to double before addition
  double(i) + d
end function
\end{lstlisting}

\section{Composite Types}
Composite types allow developers to construct structured data representations.

\subsection{Record Types}
A \texttt{record} type defines a fixed collection of named fields:
\begin{lstlisting}[caption={Record Type Declaration \& Field Access}]
type Point2D = record[ x, y : real ];
type Particle = record[ 
  id : integer; 
  pos : Point2D; 
  vel : Point2D 
];

function MoveParticle( p : Particle; dt : real returns Particle )
  p replace [ pos: record Point2D[ x: p.pos.x + p.vel.x * dt; 
                                  y: p.pos.y + p.vel.y * dt ] ]
end function
\end{lstlisting}

\subsection{Union \& Tagged Variant Types}
A \texttt{union} type represents a tagged variant type (discriminated union):
\begin{lstlisting}[caption={Union Type Declaration}]
type Shape = union[ 
  circle : record[ radius : real ]; 
  rectangle : record[ width, height : real ] 
];

function Area( s : Shape returns real )
  tagcase s
    tag circle c => 3.14159 * c.radius * c.radius
    tag rectangle r => r.width * r.height
  end tagcase
end function
\end{lstlisting}

\section{Algebraic Recursive Types}
Sisal-2026 supports recursive algebraic data structures using type references inside union definitions.

\begin{lstlisting}[caption={Algebraic Recursive List Type}]
type IntList = union[
  nil_tag  : null;
  cons_tag : record[ head : integer; tail : IntList ]
];

function SumList( L : IntList returns integer )
  tagcase L
    tag nil_tag => 0
    tag cons_tag c => c.head + SumList(c.tail)
  end tagcase
end function
\end{lstlisting}

\begin{note}[Ragged Data Structures]
Legacy Sisal 1.2 nested ragged array syntax (\texttt{array[array[T]]}) is not part of Sisal-2026. Dense multi-dimensional arrays use \texttt{array\_dv[T]}. For irregular or ragged structures (e.g. sparse trees, graphs, ragged lists), algebraic union types like \texttt{IntList} are used instead.
\end{note}
